\section*{附录}

\subsection*{附录 A：代码与数据开源说明}
本研究的全部源代码、数据处理脚本、实验配置、统计汇总程序以及论文图表生成脚本均已开源于代码仓库：\url{https://github.com/frankxucd0826-source/auv-compute-or-observe}。

\begin{enumerate}
    \item \textbf{数据集来源：}主实验海流数据取自HYCOM全球海洋再分析公开服务（产品GLBy0.08，实验93.0，10米水深水平流速，区域23°--31°N、121°--129°E，2019--2024年；重复实验区为湾流33°--41°N、75°--67°W与阿古拉斯回流区35°--43°S、14°--22°E，产品与时间范围相同；卫星地图检验所用的近实时融合测高地转流取自美国国家海洋和大气管理局海岸观测网数据服务）；第3节前期研究的叶绿素与海表温度数据取自美国国家海洋和大气管理局海岸观测网数据服务的公开卫星产品。
    \item \textbf{主要程序结构：}核心软件包 \texttt{src/oceanroute/} 包含网格最短路与包络定理梯度求解器（\texttt{solver.py}）、海流快照与粗图（\texttt{fields.py}，海区由环境变量 \texttt{OCEANROUTE\_REGION} 选择）、旧粗图误差的多年数据标定与无散误差模态（\texttt{errormodel.py}）以及任务生成、反例搜索与测量更新（\texttt{mission.py}）；数据下载、批量评估、统计汇总、随机合理场对照与湾流实验分别由 \texttt{scripts/download\_hycom.py} 与 \texttt{experiments/route\_*.py} 完成，正文插图由 \texttt{paper/figures/fig\_route\_*.py} 生成；逐条复现命令与前期研究的脚本见仓库根目录的 \texttt{README}。
\end{enumerate}

\subsection*{附录 B：人工智能工具使用说明}
本研究与论文撰写中人工智能工具的使用情况如下：
\begin{enumerate}
    \item \textbf{作者：}提出了机载计算与现场观测如何权衡的初始问题；对第一版结果提出用真实卫星数据检验、补充对照与消融、按学术论文格式写作等要求；在前期研究结果出来后指出其提升小、新意不足，并提出“粗分辨率海流图下两条航路的合理反例”这一新问题，决定转向本文研究；对全文负责。
    \item \textbf{Claude（Anthropic）：}在作者授权下完成文献检索与核对、数据下载与全部代码、实验与统计、论文插图，整理写作提纲与由实验程序生成的事实清单，并逐项核对正文中的数字与引用。
    \item \textbf{Gemini（Google）：}根据Claude整理的提纲、术语表与事实清单撰写论文文字，再由Claude核对与修改。
    \item \textbf{可追溯性记录：}作者与人工智能工具的对话记录、与Gemini的每轮写作往来以及研究笔记均保存在上述代码仓库中（涉及作者其他未发表工作的内容已隐去）。
\end{enumerate}
